\documentclass[../../main.tex]{subfiles}% 注意这里的文件路径不能用 ./main.tex ,否则用latexmk编译子文件会报错
\graphicspath{{\subfix{./image/}}{\subfix{../../image/}}} % 指定图片引用路径,后续可以直接使用图片文件名
% 如果图片存放在上级目录,则使用 ../../image/ 作为备用路径

% 例如:
% \begin{figure}[H]
% \centering
% \includegraphics[width=0.3\textwidth]{图.png}
% \caption{}
% \label{figure:图}
% \end{figure}
% 注意:上述\label{}一定要放在\caption{}之后,否则引用图片序号会只会显示??.

\begin{document}
% SOURCE p=108; printed=93; object=ch03-complex-measures-title; region=upper
\section{复测度}\label{section:ch03-complex-measures}
% SOURCE p=108; printed=93; object=ch03-complex-measure; region=upper
\begin{definition}[复测度]\label{definition:ch03-complex-measure}
可测空间 $(X,\mathscr{M})$ 上的\textbf{复测度}是映射 $\nu:\mathscr{M}\to\mathbb{C}$，满足：
\begin{enumerate}
\item $\nu(\varnothing)=0$。
\item 若 $\{E_j\}$ 为 $\mathscr{M}$ 中两两不交的集合序列，则
\begin{align*}
\nu\left(\bigcup_1^\infty E_j\right)=\sum_1^\infty\nu(E_j),
\end{align*}
且该级数绝对收敛。
\end{enumerate}
\end{definition}
\begin{remark}
特别地，复测度不允许取无穷值，所以正测度只有在有限时才是复测度。
\end{remark}
% SOURCE p=108; printed=93; object=ch03-complex-density-example; region=upper-middle
\begin{example}\label{example:ch03-complex-density}
若 $\mu$ 为正测度，$f\in L^1(\mu)$，则 $f\,\mathrm{d}\mu$ 为复测度。
\end{example}
% SOURCE p=108; printed=93; object=ch03-complex-real-imaginary; region=middle
\begin{definition}[复测度的实部与虚部]\label{definition:ch03-complex-real-imaginary}
若 $\nu$ 为复测度，用 $\nu_r,\nu_i$ 分别表示 $\nu$ 的实部和虚部。
\end{definition}
\begin{proposition}\label{proposition:ch03-complex-bounded-range}
复测度 $\nu$ 的值域是 $\mathbb{C}$ 的有界子集。
\end{proposition}
\begin{proof}
$\nu_r,\nu_i$ 为不取值 $\pm\infty$ 的符号测度，因而它们是有限的，结论随即成立。
\end{proof}
% SOURCE p=108; printed=93; object=ch03-complex-integration; region=middle
符号测度的有关概念很容易推广到复测度。例如：
\begin{definition}[关于复测度的积分]\label{definition:ch03-complex-integration}
定义 $L^1(\nu)\triangleq L^1(\nu_r)\cap L^1(\nu_i)$；对 $f\in L^1(\nu)$，令
\begin{align*}
\int f\,\mathrm{d}\nu\triangleq\int f\,\mathrm{d}\nu_r+\mathrm{i}\int f\,\mathrm{d}\nu_i.
\end{align*}
\end{definition}
% SOURCE p=108; printed=93; object=ch03-complex-singular-absolute; region=middle
\begin{definition}[复测度的奇异性与绝对连续性]\label{definition:ch03-complex-singular-absolute}
若 $\nu,\mu$ 为复测度，当对 $a,b=r,i$ 都有 $\nu_a\perp\mu_b$ 时，称 $\nu\perp\mu$。若 $\lambda$ 为正测度，当 $\nu_r\ll\lambda$ 且 $\nu_i\ll\lambda$ 时，称 $\nu\ll\lambda$。
\end{definition}
% SOURCE p=108; printed=93; object=ch03-complex-lrn-intro; region=middle
\cref{section:ch03-lebesgue-radon-nikodym}中的定理也可以推广，只需分别用于实部和虚部。特别地：
% SOURCE p=108; printed=93; object=ch03-complex-lrn; number=3.12; region=middle-lower
\begin{theorem}[Lebesgue–Radon–Nikodym 定理]\label{theorem:ch03-complex-lebesgue-radon-nikodym}
若 $\nu$ 为复测度，$\mu$ 为 $(X,\mathscr{M})$ 上的 $\sigma$-有限正测度，则存在复测度 $\lambda$ 和 $f\in L^1(\mu)$，使 $\lambda\perp\mu$ 且 $\mathrm{d}\nu=\mathrm{d}\lambda+f\,\mathrm{d}\mu$。若还有 $\lambda'\perp\mu$ 且 $\mathrm{d}\nu=\mathrm{d}\lambda'+f'\,\mathrm{d}\mu$，则 $\lambda=\lambda'$，且 $f=f'$，$\mu$-几乎处处成立。
\end{theorem}
% SOURCE p=108; printed=93; object=ch03-complex-radon-nikodym; region=lower
\begin{definition}[复测度的 Radon–Nikodym 导数]\label{definition:ch03-complex-radon-nikodym}
与前面一样，若 $\nu\ll\mu$，将\cref{theorem:ch03-complex-lebesgue-radon-nikodym}中的 $f$ 记作 $\mathrm{d}\nu/\mathrm{d}\mu$。
\end{definition}
% SOURCE p=108; printed=93; object=ch03-complex-total-variation; region=lower
\begin{definition}[复测度的全变差]\label{definition:ch03-complex-total-variation}
复测度 $\nu$ 的\textbf{全变差}是由下述性质确定的正测度 $|\nu|$：若 $\mathrm{d}\nu=f\,\mathrm{d}\mu$，其中 $\mu$ 为正测度，则 $\mathrm{d}|\nu|=|f|\,\mathrm{d}\mu$。
\end{definition}
% SOURCE p=108,109; printed=93,94; object=ch03-complex-total-variation-well-defined; region=lower/top
\begin{proposition}\label{proposition:ch03-complex-total-variation-well-defined}
上述 $|\nu|$ 的定义与 $\mu$ 和 $f$ 的选择无关，并在 $\nu$ 为符号测度时与先前的定义一致。
\end{proposition}
\begin{proof}
首先，每个 $\nu$ 都可以表示为 $f\,\mathrm{d}\mu$，其中 $\mu$ 为某个有限测度，$f\in L^1(\mu)$。事实上，可以取 $\mu=|\nu_r|+|\nu_i|$，再用\cref{theorem:ch03-complex-lebesgue-radon-nikodym}求得 $f$。其次，若 $\mathrm{d}\nu=f_1\,\mathrm{d}\mu_1=f_2\,\mathrm{d}\mu_2$，令 $\rho=\mu_1+\mu_2$。由\cref{proposition:ch03-radon-nikodym-chain}，
\begin{align*}
f_1\frac{\mathrm{d}\mu_1}{\mathrm{d}\rho}\,\mathrm{d}\rho
=\mathrm{d}\nu
=f_2\frac{\mathrm{d}\mu_2}{\mathrm{d}\rho}\,\mathrm{d}\rho,
\end{align*}
故 $f_1(\mathrm{d}\mu_1/\mathrm{d}\rho)=f_2(\mathrm{d}\mu_2/\mathrm{d}\rho)$，$\rho$-几乎处处成立。由于 $\mathrm{d}\mu_j/\mathrm{d}\rho$ 非负，得到
\begin{align*}
|f_1|\frac{\mathrm{d}\mu_1}{\mathrm{d}\rho}
=\left|f_1\frac{\mathrm{d}\mu_1}{\mathrm{d}\rho}\right|
=\left|f_2\frac{\mathrm{d}\mu_2}{\mathrm{d}\rho}\right|
=|f_2|\frac{\mathrm{d}\mu_2}{\mathrm{d}\rho}
\qquad\rho\text{-几乎处处},
\end{align*}
从而
\begin{align*}
|f_1|\,\mathrm{d}\mu_1
=|f_1|\frac{\mathrm{d}\mu_1}{\mathrm{d}\rho}\,\mathrm{d}\rho
=|f_2|\frac{\mathrm{d}\mu_2}{\mathrm{d}\rho}\,\mathrm{d}\rho
=|f_2|\,\mathrm{d}\mu_2.
\end{align*}
所以 $|\nu|$ 的定义与 $\mu,f$ 的选择无关。若 $\nu$ 为符号测度，则 $\mathrm{d}\nu=(\chi_P-\chi_N)\,\mathrm{d}|\nu|$，其中 $X=P\cup N$ 为 Hahn 分解，且 $|\chi_P-\chi_N|=1$，所以此定义与先前的定义一致。
\end{proof}
% SOURCE p=109; printed=94; object=ch03-complex-variation-properties; number=3.13; region=upper
\begin{proposition}\label{proposition:ch03-complex-variation-properties}
设 $\nu$ 为 $(X,\mathscr{M})$ 上的复测度。
\begin{enumerate}
\item\label{item:ch03-complex-variation-domination} 对所有 $E\in\mathscr{M}$，$|\nu(E)|\le|\nu|(E)$。
\item\label{item:ch03-complex-polar-density} $\nu\ll|\nu|$，且 $\mathrm{d}\nu/\mathrm{d}|\nu|$ 的绝对值 $|\nu|$-几乎处处等于 $1$。
\item\label{item:ch03-complex-integral-variation} $L^1(\nu)=L^1(|\nu|)$，并且若 $f\in L^1(\nu)$，则 $|\int f\,\mathrm{d}\nu|\le\int|f|\,\mathrm{d}|\nu|$。
\end{enumerate}
\end{proposition}
\begin{proof}
如 $|\nu|$ 的定义中那样，设 $\mathrm{d}\nu=f\,\mathrm{d}\mu$。则
\begin{align*}
|\nu(E)|=\left|\int_E f\,\mathrm{d}\mu\right|\le\int_E|f|\,\mathrm{d}\mu=|\nu|(E).
\end{align*}
这证明了\subcref{proposition:ch03-complex-variation-properties}{item:ch03-complex-variation-domination}，并说明 $\nu\ll|\nu|$。若 $g=\mathrm{d}\nu/\mathrm{d}|\nu|$，则 $f\,\mathrm{d}\mu=\mathrm{d}\nu=g\,\mathrm{d}|\nu|=g|f|\,\mathrm{d}\mu$，故 $g|f|=f$，$\mu$-几乎处处成立，从而也 $|\nu|$-几乎处处成立。显然 $|f|>0$，$|\nu|$-几乎处处成立，所以 $|g|=1$，$|\nu|$-几乎处处成立。\subcref{proposition:ch03-complex-variation-properties}{item:ch03-complex-integral-variation}留给读者证明（\cref{exercise:ch03-complex-integral-variation}）。
\end{proof}
% SOURCE p=109; printed=94; object=ch03-complex-variation-triangle; number=3.14; region=middle
\begin{proposition}\label{proposition:ch03-complex-variation-triangle}
若 $\nu_1,\nu_2$ 为 $(X,\mathscr{M})$ 上的复测度，则 $|\nu_1+\nu_2|\le|\nu_1|+|\nu_2|$。
\end{proposition}
\begin{proof}
由\cref{proposition:ch03-common-dominating-measure}，对 $j=1,2$，可以用同一个 $\mu$ 写成 $\mathrm{d}\nu_j=f_j\,\mathrm{d}\mu$。于是
\begin{align*}
\mathrm{d}|\nu_1+\nu_2|
=|f_1+f_2|\,\mathrm{d}\mu
\le|f_1|\,\mathrm{d}\mu+|f_2|\,\mathrm{d}\mu
=\mathrm{d}|\nu_1|+\mathrm{d}|\nu_2|.
\end{align*}
\end{proof}
% SOURCE p=109; printed=94; object=ch03-complex-exercises-title; region=middle
\subsection*{习题}
% SOURCE p=109; printed=94; object=ch03-ex18; number=18; region=middle
\begin{exercise}\label{exercise:ch03-complex-integral-variation}
证明\subcref{proposition:ch03-complex-variation-properties}{item:ch03-complex-integral-variation}。
\end{exercise}
% SOURCE p=109; printed=94; object=ch03-ex19; number=19; region=middle-lower
\begin{exercise}\label{exercise:ch03-complex-singular-absolute}
若 $\nu,\mu$ 为复测度，$\lambda$ 为正测度，则 $\nu\perp\mu$ 当且仅当 $|\nu|\perp|\mu|$，而 $\nu\ll\lambda$ 当且仅当 $|\nu|\ll\lambda$。
\end{exercise}
% SOURCE p=109; printed=94; object=ch03-ex20; number=20; region=lower
\begin{exercise}\label{exercise:ch03-complex-positive-criterion}
若 $\nu$ 为 $(X,\mathscr{M})$ 上的复测度，且 $\nu(X)=|\nu|(X)$，则 $\nu=|\nu|$。
\end{exercise}
% SOURCE p=109,110; printed=94,95; object=ch03-ex21; number=21; region=lower/top
\begin{exercise}\label{exercise:ch03-complex-variation-suprema}
设 $\nu$ 为 $(X,\mathscr{M})$ 上的复测度。若 $E\in\mathscr{M}$，定义
\begin{align*}
\mu_1(E)&=\sup\left\{\sum_1^n|\nu(E_j)|:n\in\mathbb{N},\ E_1,\ldots,E_n\text{ 两两不交},\ E=\bigcup_1^nE_j\right\},\\
\mu_2(E)&=\sup\left\{\sum_1^\infty|\nu(E_j)|:E_1,E_2,\ldots\text{ 两两不交},\ E=\bigcup_1^\infty E_j\right\},\\
\mu_3(E)&=\sup\left\{\left|\int_E f\,\mathrm{d}\nu\right|:|f|\le1\right\}.
\end{align*}
则 $\mu_1=\mu_2=\mu_3=|\nu|$。（先证明 $\mu_1\le\mu_2\le\mu_3$。为证明 $\mu_3=|\nu|$，令 $f=\overline{\mathrm{d}\nu/\mathrm{d}|\nu|}$，并应用\cref{proposition:ch03-complex-variation-properties}。为证明 $\mu_3\le\mu_1$，用简单函数逼近 $f$。）
\end{exercise}
\end{document}
